\sheettitle
  {Feuille de travail}
  {$L^p$ spaces and Radon--Nikodym representations}

Unless otherwise stated, $(E,\mathcal A,\mu)$ is a measure space. Formal statements, estimates, and almost-everywhere bases must retain their hypotheses.

\section*{Exercise 0 — The probability-space $L^p$ structure}

Let $(\Omega,\mathcal F,\mathbb P)$ be a probability space.
\begin{enumerate}
  \item Define $L^p(\mathbb P)$ as a quotient space. State the equivalence relation and explain why the ambient measure is part of the object.
  \item State
  \[
  L^\infty(\mathbb P)
  \subseteq
  L^2(\mathbb P)
  \subseteq
  L^1(\mathbb P)
  \]
  together with the norm inequalities. Identify exactly where $\mathbb P(\Omega)=1$ enters. Do not re-prove the embedding in this part.
  \item Translate $X\in L^p(\mathbb P)$ into a moment statement.
  \item State the estimates which justify
  \[
  L^1\times L^\infty\longrightarrow L^1,
  \qquad
  L^2\times L^2\longrightarrow L^1.
  \]
  Name the theorem used at each endpoint.
  \item Assign $L^1$, $L^\infty$, and $L^2$ respectively to finite expectation and RN numerators, bounded tests, and conditional second moments.
\end{enumerate}

\emph{Time limit: 15 minutes. This is an object-and-role reconstruction.}

\section*{Exercise 1 — The finite-measure embedding}

Assume $\mu(E)<\infty$ and let $1\le p<q\le\infty$.
\begin{enumerate}
  \item For $q<\infty$, prove
  \[
  \|f\|_p
  \le
  \mu(E)^{1/p-1/q}\|f\|_q.
  \]
  Display the two conjugate Hölder exponents and identify the precise use of $\mu(E)<\infty$.
  \item Treat $q=\infty$ and specialize the result to a probability measure. State the resulting hierarchy between $L^\infty$, $L^2$, and $L^1$.
\end{enumerate}

\emph{Time limit: 20 minutes. Sharpness and long infinite-measure counterexamples are deferred.}

\section*{Exercise 2 — From indicators to bounded tests}

Let $\mathcal G\subseteq\mathcal A$ be a sub-$\sigma$-field and let $f,g\in L^1(\mu)$ satisfy
\[
\int_Af\,d\mu
=
\int_Ag\,d\mu
\qquad(A\in\mathcal G).
\]
\begin{enumerate}
  \item Prove that $YZ\in L^1(\mu)$ and
  \[
  \|YZ\|_1\le\|Y\|_1\|Z\|_\infty
  \]
  whenever $Y\in L^1(\mu)$ and $Z\in L^\infty(\mu)$.
  \item Deduce the testing identity for every bounded $\mathcal G$-measurable simple function.
  \item Extend the identity to every bounded $\mathcal G$-measurable function. State the dominating integrable function.
  \item If $f$ and $g$ are $\mathcal G$-measurable, prove that $f=g$ $\mu$-almost everywhere.
\end{enumerate}

\emph{Time limit: 20 minutes. The final passage must name dominated convergence.}

\section*{Exercise 3 — Density measures and change of measure}

\begin{enumerate}
  \item Let $f\in L^1(\mu)$. Prove that
  \[
  \nu_f(A)=\int_Af\,d\mu
  \]
  defines a finite signed measure and that $\nu_f\ll\mu$.
  \item If $f,g\in L^1(\mu)$ and $\nu_f=\nu_g$, prove that $f=g$ $\mu$-almost everywhere.
  \item Let $\nu=h\mu$ with $h\ge0$. Prove
  \[
  \int_E\varphi\,d\nu
  =
  \int_E\varphi h\,d\mu
  \]
  first for nonnegative $\varphi$, then for signed $\varphi$ in the correct integrability regime. Deduce the criterion for $\varphi\in L^1(\nu)$.
\end{enumerate}

\emph{Time limit: 20 minutes. Forward density generation must not be confused with an RN converse.}

\section*{Exercise 4 — Four uses of Radon--Nikodym}

Complete the following table.

\begin{center}
\begin{tabular}{p{32mm}p{34mm}p{31mm}p{30mm}p{21mm}}
\textbf{Context} & \textbf{Numerator} & \textbf{Reference / condition} & \textbf{Domain and a.e. basis} & \textbf{Date}\\
\hline
density generated by $f\in L^1(\mu)$ & & & & 14 July\\[2.0em]
conditioning on $B$, $\mathbb P(B)>0$ & & & & 19 July\\[2.0em]
density of the law of an $\mathbb R^d$-valued $X$, assuming $\mu_X\ll\lambda_d$ & & & & 24 July\\[2.0em]
conditional expectation of $Y\in L^1(\mathbb P)$ on a sub-$\sigma$-field $\mathcal G\subseteq\mathcal F$ & & & & 26 July
\end{tabular}
\end{center}

Then answer:
\begin{enumerate}
  \item Which rows are explicit forward density constructions, and which require an RN converse?
  \item Why is $\mu_X=X_\#\mathbb P$ already a law but not yet a density?
  \item How is the positive RN theorem legitimately extended to the signed conditional-expectation numerator?
  \item Relative to which measure is each representative unique?
\end{enumerate}

\emph{Time limit: 10 minutes. Do not reconstruct a Hahn-first proof.}

\section*{Closed-book reconstruction}

Complete and classify the arrows:
\[
\boxed{Y\in L^q(\mathbb P),\ q>1}
\Longrightarrow
\boxed{Y\in L^1(\mathbb P)}
\Longrightarrow
\boxed{Y\mathbb P\text{ is finite signed}}
\Longrightarrow
\boxed{\text{RN representation after a reference measure is chosen}}.
\]
The final arrow is legitimate only after the numerator, reference measure, absolute continuity, and RN regime have been identified.

\section*{Further work}

The full Hahn local-density fragment, maximal represented mass, residue elimination, $\sigma$-finite gluing, long infinite-measure counterexamples, sharpness, and full $L^2$ projection theory are outside this sheet.
